\documentclass[10pt,a4paper]{article}

% =========================
% Packages
% =========================
\usepackage[utf8]{inputenc}
\usepackage[T1]{fontenc}
\usepackage[margin=0.85in]{geometry}
\usepackage{times}
\usepackage{xcolor}
\usepackage{titlesec}
\usepackage{setspace}
\usepackage{enumitem}
\usepackage{float}
\usepackage{booktabs}
\usepackage{multirow}
\usepackage{array}
\usepackage{tabularx}
\usepackage{fancyhdr}
\usepackage[hidelinks]{hyperref}

% =========================
% Formatting (Same Template Style)
% =========================
\definecolor{amritaMaroon}{HTML}{A4123F}
\titleformat{\section}{\Large\bfseries\color{amritaMaroon}}{\thesection}{1em}{}
\titleformat{\subsection}{\large\bfseries}{\thesubsection}{1em}{}
\titlespacing*{\section}{0pt}{6pt}{2pt}
\titlespacing*{\subsection}{0pt}{4pt}{1pt}
\setstretch{1.0}
\setlist[itemize]{noitemsep,topsep=1pt,leftmargin=5mm}

\pagestyle{fancy}
\fancyhf{}
\fancyhead[L]{Phase 1 Findings Summary}
\fancyhead[R]{\thepage}
\renewcommand{\headrulewidth}{0.4pt}

\newcolumntype{L}{>{\raggedright\arraybackslash}X}

\begin{document}

% =========================
% Cover Block
% =========================
\begin{center}
{\Large\textbf{\textcolor{amritaMaroon}{AMRITA VISHWA VIDYAPEETHAM}}}\\
{\large\textbf{School of Computing}}\\
{\large\textbf{Department of Computer Science and Engineering}}\\
\vspace{0.2cm}
{\LARGE\textbf{Phase 1 Summary Report}}\\
\vspace{0.1cm}
{\large\textbf{Summary of Findings and Completion Status}}\\
\vspace{0.1cm}
23CSE399 -- Structured Two-Page Note\\
\today
\end{center}

\vspace{0.1cm}
\hrule
\vspace{0.08cm}

\begin{center}
\renewcommand{\arraystretch}{1.25}
\setlength{\tabcolsep}{3pt}
\begin{tabular}{|l|p{4.2cm}|p{4cm}|p{2.6cm}|}
\hline
\multicolumn{2}{|c|}{\textbf{Team Members}} & \textbf{Guide} & \textbf{Project Type} \\
\hline
CB.SC.U4CSE23018 & Rohith Kumar D &
\multirow{5}{=}{\raggedright \small \textbf{Dr. Venkataraman D} \\ \textit{Associate Professor, Dept. of CSE}} &
\multirow{5}{=}{\raggedright \small \textbf{Type 1} \\ Research} \\ \cline{1-2}
CB.SC.U4CSE23055 & T Venkataramana & & \\ \cline{1-2}
CB.SC.U4CSE23058 & VAB Jashwanth Reddy & & \\ \cline{1-2}
CB.SC.U4CSE23060 & C Kalyan Kumar Reddy & & \\ \cline{1-2}
CB.SC.U4CSE23446 & Hemanth Sholingaram & & \\
\hline
\end{tabular}
\end{center}

\section{Executive Summary of Phase 1 Findings}
\small
Phase 1 established the research direction, novelty scope, and system ambition for an edge-based AI speech therapy framework focused on children with ASD in regional language settings. The report and panel review highlighted both strengths and corrections needed for realistic deployment.

\subsection{Major Findings}
\begin{itemize}
\item \textbf{Strong Foundation:} The project clearly justified clinical need, regional language relevance, and therapist-access constraints.
\item \textbf{Comprehensive Literature Base:} A broad survey (25 papers) covered pronunciation metrics, multimodal cues, adaptive learning, and privacy frameworks.
\item \textbf{Core Technical Direction Validated:} Syllable-level scoring, VTLN normalization, and longitudinal tracking were accepted as meaningful.
\item \textbf{Feasibility Gap Identified:} The original architecture was too heavy for immediate edge deployment due to concurrent high-compute modules.
\item \textbf{Adaptivity Timing Concern:} RL was considered useful but unsafe for immediate use without sufficient interaction data.
\end{itemize}

\subsection{Panel-Driven Corrections Adopted}
\begin{itemize}
\item Shift from full-stack immediate deployment to staged integration.
\item Replace computationally expensive privacy approach with practical local-first federated synchronization.
\item Prioritize deterministic scoring pipeline before advanced adaptive modules.
\item Improve architecture clarity with explicit edge vs sync boundary.
\end{itemize}

\begin{table}[htbp]
\centering
\caption{Phase 1 Findings to Engineering Decision Mapping}
\renewcommand{\arraystretch}{1.2}
\begin{tabularx}{\textwidth}{|p{4cm}|L|}
\hline
\textbf{Phase 1 Finding} & \textbf{Decision Taken for Next Stage} \\
\hline
Stack too heavy for edge at once & Introduce phased module rollout \\
\hline
Latency claim not realistic under full load & Optimize core pipeline first \\
\hline
FHE overhead impractical for target hardware & Use practical federated update flow \\
\hline
Immediate RL is data-unsafe & Keep rule-based control before RL \\
\hline
Tracking layer over-complex in early stage & Use structured lightweight progress tracking \\
\hline
\end{tabularx}
\end{table}

\vspace{0.1cm}

\section{Completion of Phase 1: Structured Status}
\subsection{Completed Outcomes}
\begin{itemize}
\item Problem definition and research motivation were fully established.
\item Literature review and cumulative research-gap identification were completed.
\item Originality claims were framed and defended with comparative justification.
\item Baseline architecture and module decomposition were documented.
\item Initial evaluation metrics and expected outcomes were defined.
\item Dataset strategy and proxy rationale were identified.
\item Major risks (latency, privacy overhead, RL convergence) were documented.
\end{itemize}

\subsection{What Phase 1 Did Not Finalize (By Design)}
\begin{itemize}
\item End-to-end edge deployment benchmarking was not fully finalized.
\item RL policy deployment was intentionally deferred pending data readiness.
\item Primary target-population corpus collection remains a planned controlled phase.
\item Clinical pilot validation was scoped, not completed, in Phase 1.
\end{itemize}

\subsection{Completion Matrix}
\begin{table}[htbp]
\centering
\renewcommand{\arraystretch}{1.12}
\setlength{\tabcolsep}{3.5pt}
\begin{tabularx}{\textwidth}{|p{5.1cm}|p{2.2cm}|L|}
\hline
\textbf{Phase 1 Component} & \textbf{Status} & \textbf{Remarks} \\
\hline
Problem formulation and objective framing & Completed & Clear and clinically motivated \\
\hline
Literature survey and gap synthesis & Completed & Sufficient for Phase 2 transition \\
\hline
Architecture and novelty articulation & Completed with revisions & Feasibility refinements incorporated \\
\hline
Core metric definition (GOP/VSA/latency intent) & Completed & Ready for Phase 2 validation \\
\hline
Full deployment and pilot evidence & Pending next phase & Correctly deferred to avoid premature claims \\
\hline
\end{tabularx}
\end{table}

\enlargethispage{3\baselineskip}

\subsection{Final Phase 1 Completion Statement}
Phase 1 is considered \textbf{successfully completed} as a research-definition and feasibility-discovery stage. It delivered a validated problem statement, strong technical foundation, and precise correction map for Phase 2 execution. The transition quality is high because critical constraints were identified early and translated into actionable engineering decisions.

\begin{center}
\footnotesize \textbf{End of Two-Page Summary}
\end{center}

\end{document}
